\section{The Four-State Orientation Cell}
\label{app:v4}
Identify a fiber with $V_4\cong\F_2^2$ and write $a=(1,0)$, $b=(0,1)$. The multiplication table is
\begin{center}
\begin{tabular}{c|cccc}
 & $1$ & $a$ & $b$ & $ab$\\\hline
$1$ & $1$ & $a$ & $b$ & $ab$\\
$a$ & $a$ & $1$ & $ab$ & $b$\\
$b$ & $b$ & $ab$ & $1$ & $a$\\
$ab$ & $ab$ & $b$ & $a$ & $1$
\end{tabular}
\end{center}
Translation by each nonzero element is a fixed-point-free involution. Its two unordered pairs are
\[
\begin{array}{c|c}
a & \{1,a\},\ \{b,ab\}\\
b & \{1,b\},\ \{a,ab\}\\
ab & \{1,ab\},\ \{a,b\}.
\end{array}
\]
Together these are the three perfect matchings partitioning the edges of an abstract four-vertex complete graph. In the reconstructed carrier, the members of a quotient fiber have no carrier edges between them; the package checks this explicitly. The matching interpretation describes deck relations, not a literal induced $K_4$ subgraph.

Any automorphism of $V_4$ permutes its three nonidentity elements, and every such permutation is an automorphism because the product of any two distinct nonidentity elements is the third. Thus $\Aut(V_4)\cong S_3$. This abstract triality does not prove that every permutation extends to a symmetry of a chosen event grammar or preserves a chosen phase, source, or release rule. Those extensions are additional mathematical questions.

The involution $a$ is an operation on a fiber, not a vertex representing a cubic centroid. Likewise, the three matchings are not three demonstrated Euclidean spatial axes. A geometric presentation may illustrate them, but its interpretation must not change the group-theoretic type of the objects.
